% Generated by tools/docs/gen_fact_tables.py from rtl/soc/apb_cdc_bridge.sv and docs/manual/common/facts/cdc.yaml. Do not edit.
{\small
\begin{longtable}{@{}>{\raggedright\arraybackslash}p{31mm}>{\raggedright\arraybackslash}p{9mm}>{\raggedright\arraybackslash}p{19mm}>{\raggedright\arraybackslash}p{22mm}>{\raggedright\arraybackslash}p{58mm}@{}}
\caption{Signals of \code{apb_cdc_bridge}}\label{tab:tm:mod:apb-cdc-bridge:signals}\\
\toprule
Signal & Dir & Width & Clock domain & Description\\
\midrule
\endfirsthead
\multicolumn{5}{@{}l}{\emph{Table \thetable\ continued}}\\
\toprule
Signal & Dir & Width & Clock domain & Description\\
\midrule
\endhead
\bottomrule
\endlastfoot
\code{scanmode_i} & in & 1 & quasi-static & Scan mode. It goes to both internal reset synchronisers. Tie low without scan.\\
\code{scan_rst_ni} & in & 1 & async & Active-low scan reset. It goes to both internal reset synchronisers. Tie high without scan.\\
\code{s_clk_i} & in & 1 & clock & Clock of the source face.\\
\code{s_rst_n_i} & in & 1 & async & Active-low reset of the source face. The bridge synchronises it inside. It also resets the destination side.\\
\code{s_psel_i} & in & 1 & \code{s_clk_i} & APB select from the upstream master.\\
\code{s_penable_i} & in & 1 & \code{s_clk_i} & APB enable from the upstream master.\\
\code{s_pwrite_i} & in & 1 & \code{s_clk_i} & APB write flag. The bridge captures it in the SETUP cycle.\\
\code{s_paddr_i} & in & \code{ADDR_W} (12) & \code{s_clk_i} & APB byte address. The bridge captures it in the SETUP cycle.\\
\code{s_pwdata_i} & in & 32 & \code{s_clk_i} & APB write data. The bridge captures it in the SETUP cycle.\\
\code{s_pstrb_i} & in & 4 & \code{s_clk_i} & APB byte strobes. The bridge captures them in the SETUP cycle.\\
\code{s_prdata_o} & out & 32 & \code{s_clk_i} & APB read data. It holds the data of the last completed transfer. It is zero after an error completion that the bridge forces.\\
\code{s_pready_o} & out & 1 & \code{s_clk_i} & APB ready. It is high for one cycle at the end of each transfer.\\
\code{s_pslverr_o} & out & 1 & \code{s_clk_i} & APB error. It holds the error flag of the last completed transfer. It is high after an error completion that the bridge forces.\\
\code{m_clk_i} & in & 1 & clock & Clock of the destination face.\\
\code{m_rst_n_i} & in & 1 & async & Active-low reset of the destination face. The bridge synchronises it inside. It also crosses into the \code{s_clk_i} domain as a status.\\
\code{m_psel_o} & out & 1 & \code{m_clk_i} & APB select to the downstream slave. High in the SETUP and ACCESS phases.\\
\code{m_penable_o} & out & 1 & \code{m_clk_i} & APB enable to the downstream slave. High in the ACCESS phase.\\
\code{m_pwrite_o} & out & 1 & \code{m_clk_i} & APB write flag of the captured transfer. It holds its value when the bridge is idle.\\
\code{m_paddr_o} & out & \code{ADDR_W} (12) & \code{m_clk_i} & APB byte address of the captured transfer. It holds its value when the bridge is idle.\\
\code{m_pwdata_o} & out & 32 & \code{m_clk_i} & APB write data of the captured transfer. It holds its value when the bridge is idle.\\
\code{m_pstrb_o} & out & 4 & \code{m_clk_i} & APB byte strobes of the captured transfer. They hold their value when the bridge is idle.\\
\code{m_prdata_i} & in & 32 & \code{m_clk_i} & APB read data from the downstream slave.\\
\code{m_pready_i} & in & 1 & \code{m_clk_i} & APB ready from the downstream slave.\\
\code{m_pslverr_i} & in & 1 & \code{m_clk_i} & APB error from the downstream slave.\\
\end{longtable}
}
